% SEGMENT OLP-0037-S01
% Part:sets-functions-relations
% Chapter: sets
% Section: schroder-bernstein

\documentclass[../../../include/open-logic-section]{subfiles}

\begin{document}

\olfileid{sfr}{siz}{sb}

\olsection{அளவு என்ற கருத்தும் ஷ்ரோடர்--பெர்ன்ஸ்டைன் தேற்றமும்}

\begin{explain}
ஓர் உள்ளுணர்வுக் கருத்தை எடுத்துக்கொள்வோம்:
$A$ என்பது $B$ ஐ விடப் பெரிதல்ல என்றும்,
$B$ என்பது $A$ ஐ விடப் பெரிதல்ல என்றும் இருந்தால்,
$A$, $B$ ஆகியவை எண்ணொத்தவை.
உண்மையாகச் சொன்னால், இந்தக் கருத்து \emph{தவறாக} இருந்தால்,
நாம் வரையறுத்த எண்ணொத்த தன்மைக்கும் கணங்களின் ``அளவுகளை''
ஒப்பிடுவதற்கும் தொடர்புண்டு என்ற எண்ணத்தை நியாயப்படுத்துவது
மிகவும் கடினமாகிவிடும்!{}
நல்வாய்ப்பாக, அந்த உள்ளுணர்வுக் கருத்து சரியானது.
ஷ்ரோடர்--பெர்ன்ஸ்டைன் தேற்றம் இதனை நியாயப்படுத்துகிறது.
\end{explain}

% SEGMENT OLP-0037-S02
\begin{thm}[ஷ்ரோடர்--பெர்ன்ஸ்டைன்]
	\ollabel{thm:schroder-bernstein}
	$\cardle{A}{B}$ மற்றும் $\cardle{B}{A}$ எனில்,
	$\cardeq{A}{B}$.
\end{thm}

% SEGMENT OLP-0037-S03
\begin{explain}
வேறு சொற்களில், $A$ இலிருந்து $B$ க்கு !!a{injection} உம்,
$B$ இலிருந்து $A$ க்கு !!a{injection} உம் இருந்தால்,
$A$ இலிருந்து $B$ க்கு !!a{bijection} இருக்கும்.

ஆயினும், இந்த முடிவை நிறுவுவது உண்மையில் மிகவும் \emph{கடினம்}.
காண்டோர் இந்த முடிவைக் கூறியிருந்தாலும், பிறரே அதனை நிறுவினர்.
\footnote{வரலாற்றைப் பற்றி மேலும் அறிய, எடுத்துக்காட்டாக,
\citet[pp.~165--6]{Potter2004} ஐப் பார்க்கவும்.}
\oliflabeldef{sfr:cardinals:card-sb:sec}{ஷ்ரோடர்--பெர்ன்ஸ்டைன்
தேற்றத்தை \emph{நிறுவும்} நிலையில் நாம்
\olref[sfr][cardinals][card-sb]{sec} இல் மட்டுமே இருப்போம்.}{}%
இப்போதைக்கு, அதனை நம்பி ஏற்றுக்கொள்ளலாம் (ஏற்றுக்கொள்ளவும் வேண்டும்).

நல்வாய்ப்பாக, ஷ்ரோடர்--பெர்ன்ஸ்டைன் தேற்றம் \emph{சரியானது};
மேலும் நாம் வரையறுத்த $\cardeq{A}{B}$, $\cardle{A}{B}$ ஆகிய
தொடர்புகளுக்கு ``அளவுடன்'' ஏதோ தொடர்புண்டு என்று கருதுவதை
அது உறுதிப்படுத்துகிறது.
அதோடு, ஷ்ரோடர்--பெர்ன்ஸ்டைன் தேற்றம் மிகவும் \emph{பயனுள்ளது}.
எண்ணொத்த இரண்டு கணங்களுக்கு இடையில் !!a{bijection} ஒன்றைச்
சிந்தித்துக் கண்டறிவது கடினமாக இருக்கலாம்.
ஷ்ரோடர்--பெர்ன்ஸ்டைன் தேற்றம் ஒப்பீட்டை இரண்டு திசைகளாகப் பிரிக்கிறது:
முதல் கணத்திலிருந்து இரண்டாவது கணத்திற்குச் செல்லும் !!a{injection}
ஒன்றையும், மறுதிசையில் அவ்வாறே செல்லும் ஒன்றையும் மட்டும் சிந்தித்தால் போதும்.
\end{explain}
\end{document}
